\documentclass[11pt]{article}
\usepackage[utf8]{inputenc}
\usepackage{amsmath}
\usepackage{geometry}
\geometry{margin=2cm}

\title{Notes de SVT - Biologie moléculaire}
\date{}

\begin{document}

\maketitle



\section*{Intro to SV}

\section*{19 février}

L'être humain $\rightarrow$ 70\% de la masse corporelle = eau ; 90\% de la masse du cerveau = eau\\
$\rightarrow$ \textbf{LA VIE REQUIERT DE L'EAU}

\vspace{0.3cm}

Ca tombe bien : l'eau est majoritairement composée d'atomes reliés par des liaisons chimiques covalentes, des liaisons solides et donc fortes contrairement aux liaisons peu solides, donc faibles, les ponts hydrogènes. Comment tester ? On mesure avec l'énergie et la force nécessaire pour les casser.

L'eau s'évapore lorsque les liaisons, les ponts hydrogène sont détruits par l'énergie calorifique du chauffage.

PRATIQUEMENT 10X plus d'énergie pour casser une liaison covalente que un pont hydrogène.

\vspace{0.3cm}

\textbf{Liaisons covalentes :} c'est un partage des électrons.

\textbf{Polaires et non polaires :} polaire s'il y a “injustice du partage des électrons”, on regarde dépendamment à l'électronégativité, la force à attirer les électrons. Si elle est égale alors les électrons sont parfaitement partagés.

Une molécule est polaire si toutes les liaisons sont polaires, elle est non polaire si toutes les liaisons sont non polaires, et finalement, s'il y a des liaisons polaires et non polaires, alors la molécule est amphiphile.

\section*{20 février}

Est-ce que l'air est hydrophobe ou hydrophile ? L'air est non polaire.\\
$\rightarrow$ bien que les liaisons du CO\textsubscript{2} sont polaires, on le considère comme un non polaire du fait de la symétrie vectorielle.

\vspace{0.3cm}

L'eau dans l'état liquide forme des ponts hydrogènes avec l'eau et si une molécule qui veut en former aussi arrive alors elle est acceptée, dans le cas contraire, elle est rejetée.

Elle fait en moyenne 3 PH à l'état liquide, 4 à l'état solide et ces ponts peuvent bouger et changer à l'état liquide.

\vspace{0.3cm}

Le savon : la molécule de savon est amphiphile, séparée entre une queue non polaire et une tête polaire COOH.

Par conséquent, la formation d'une bulle de savon demande une double couche de molécules de savon, l'eau est prise entre les têtes des molécules qui forment donc deux couches.

C'est le principe inverse pour les membranes cellulaires : à l'inverse on va avoir une double couche sans eau à l'intérieur et donc de l'eau à l'extérieur. Pour permettre le contact et l'échange moléculaire avec l'environnement externe à la cellule, la création de canaux ioniques permet de faire passer des molécules spécifiques. Un canal est spécifique à une molécule. Ici c'est une bicouche de phospholipides.

\vspace{0.3cm}

Quand la liaison covalente devient encore plus inégale, on parle alors de liaison ionique.

\vspace{0.3cm}

\textbf{L'ADN :}

Backbone $\rightarrow$ les deux côtés, les deux barres de l'hélice de l'ADN.

Chaque côté est fait de liaisons covalentes mais entre les deux brins, ce sont des ponts hydrogène.

Bouillir à 65°C sépare les 2 brins : c'est la dénaturation de l'ADN. Quand on reréduit la température, les brins se remettent automatiquement ensemble et les bons avec les bons.

Les deux brins sont antiparallèles 5' 3' : on va toujours dans le sens de 5' VERS le 3'.

\section*{23 février}

\textbf{ADN}

Cellule procaryote comme \textit{E. coli} : ADN circulaire.

Cellule eucaryote : noyau de 10 micromètres de diamètre.

L'ADN a toujours des phosphates et est donc chargé négativement. Chez l'humain : 46 filaments, 46 chromosomes.

\vspace{0.3cm}

\textbf{État général :}

Double hélice nue $\rightarrow$ 2 nm d'épaisseur\\
Billes sur une corde $\rightarrow$ nucléofilaments $\rightarrow$ 10 nm d'épaisseur\\
Nucléosomes : protéines (histones) chargées positivement autour desquelles s'enroule l'ADN\\
ADN condensé $\rightarrow$ 30 nm d'épaisseur

3 nucléosomes à chaque fois + 8 ?

On ne peut plus lire l'ADN à ce stade.

Puis si c'est encore plus compacté, on atteint l'état de chromosome à 1 ou 2 chromatides (2 si avant la division cellulaire), taille supérieure à 1 micromètre (environ 1.4 µm).

\vspace{0.3cm}

Comment fait-on pour avoir de l'ADN nu ? On prend des bactéries, on utilise des enzymes qui enlèvent les protéines.

\vspace{0.3cm}

Les bactéries se divisent pour se dupliquer, nous on fait la mitose :

G1 (growth 1) $\rightarrow$ S (duplication de l'ADN, synthèse) $\rightarrow$ G2 (wait) [interphase] $\rightarrow$ mitose (prophase, métaphase, anaphase, télophase)

CYCLE CELLULAIRE :

[]

\textbf{Déf : Turn over :} utilisé en génétique pour décrire la dynamique du génome ou la traduction.

\section*{26 février}

\textbf{Polymère}

Si le polymère est linéaire, il a 2 extrémités. Si elles sont différentes, alors il a un sens.

ATCG $\rightarrow$ adénine, thymine, cytosine, guanine.

Polynucléotide $\rightarrow$ chaîne d'acides nucléiques.

5' PO$_4^-$ $\quad$ 3' -OH

Les acides aminés ont un acide et une amine aux deux extrémités.

Monomère = nucléotide.

\vspace{0.3cm}

Base\\
\hspace{1cm} |\\
Sucre $\rightarrow$ désoxyribose\\
\hspace{1cm} |\\
P

Le sucre est accepté par le carbone 3 du phosphate et le carbone 4 ferme le cycle du sucre.

\vspace{0.3cm}

\textbf{Copie de l'ADN}

Coupe en 2 et chaque brin sert de modèle (template).  
On ne peut copier que si on a séparé.

(ADN humain $\rightarrow$ 1 m, ADN bactérien $\rightarrow$ 1 mm)

La réplication est bidirectionnelle.

Fourche de réplication : endroit où la séparation se fait et où on ajoute les nouveaux éléments.

Réplisome : ensemble des processus mis en œuvre pour répliquer l'ADN.

Temps de duplication : environ 20 min pour \textit{E. coli} à 37°C.

5' est l'extrémité morte : on ne lui ajoute rien, on ajoute toujours au 3'.

ADN polymérase $\rightarrow$ duplique.

Nécessite une amorce de 10 bases.

Cette amorce est de l'ARN, ensuite remplacée par de l'ADN.

Primase : enzyme qui fabrique l'amorce (ARN polymérase).

Dans le sens inverse : fragments d'Okazaki (~1000 bases).

\section*{27 février}

ADN polymérase ne va pas au 5' de l'amorce, se pose au 3'.

L'addition ne se fait que du côté 3', donc on va sur le brin original de 3' vers 5' (sur l'original mais on travaille sur l'amorce).

ADN polymérase III.

Fragments d'Okazaki : ~1000 bases.

La primase refait une amorce dès qu'il y a 1000 bases.

\textit{E. coli} : 4 milliards de bases $\approx$ 1 mm de périmètre.

Origine de réplication : OriC (chez les bactéries).

Sur OriC : tensions et rotations cassent les ponts hydrogène et séparent les brins.

Puis ADN hélicase continue la séparation.

Single-stranded DNA-binding proteins maintiennent les brins + ADN gyrase aide.

Ensuite amorces (ARN ~10 nucléotides), puis ADN polymérase.

Dans l'autre sens : fragments d'Okazaki.

ADN polymérase I remplace l'ARN par de l'ADN.

ADN ligase relie les fragments en formant la liaison covalente finale.




\section*{PCR (Polymerase Chain Reaction)}

Fait une amplification exponentielle de l'ADN.\\
1 milliard de copies après 30 cycles.\\
On utilise un tout petit volume initial : 25 µL.

\vspace{0.3cm}

\textbf{On a besoin :}
\begin{itemize}
    \item L'ADN original
    \item Une ADN polymérase : Taq ADN polymérase
    \item Des monomères
    \item Des amorces (ADN commandées sur mesure)
\end{itemize}

On met le tube dans un thermocycleur.

\vspace{0.3cm}

\textbf{Cycle :}
\begin{itemize}
    \item 90°C $\rightarrow$ dénaturation : les ponts hydrogène sont brisés et les deux brins se séparent
    \item Annealing (48 à 72°C) $\rightarrow$ les amorces se lient aux séquences complémentaires
    \item 68 à 72°C $\rightarrow$ extension : l'ADN polymérase allonge les chaînes à partir des amorces
\end{itemize}

\section*{Réplication de l'ADN}

Les DNA polymérases ont des problèmes avec les séquences répétitives en tandem (microsatellites, ex : GATA répété 12 fois = 48 bases).\\
Les microsatellites sont aussi appelés STR.

\vspace{0.3cm}

Erreurs possibles :
\begin{itemize}
    \item Substitution d'une base : environ $10^{-9}$
    \item Insertion/délétion (ex : GATA ajouté ou supprimé) : environ $10^{-3}$ (beaucoup plus fréquent)
\end{itemize}

\vspace{0.3cm}

\textbf{SNP :} Single Nucleotide Polymorphism\\
Différences :
\begin{itemize}
    \item entre humains : ~1 base sur 1000
    \item humain / chimpanzé : ~1 base sur 100
\end{itemize}

Dans l'ADN, certaines lettres ne servent à rien.

\section*{Électrophorèse}

On utilise des tubes. On observe les molécules par fluorescence (laser).\\
Type : électrophorèse capillaire / gel.

L'ADN passe dans des pores :
\begin{itemize}
    \item petites molécules : rapides
    \item grandes molécules : plus lentes
\end{itemize}

\section*{9 mars : Méiose}

Ovaires et testicules = gonades → seuls lieux de la méiose.\\
Gamètes = cellules reproductrices (germinales), haploïdes (1 chromosome de chaque paire).

\vspace{0.3cm}

Non-disjonction en méiose I.\\
Non-disjonction en méiose II (cas C).

\vspace{0.3cm}

Prophase I : les chromosomes homologues s'apparient et effectuent des échanges (crossing-over).

\vspace{0.3cm}

Méiose chez l'homme :
\begin{itemize}
    \item 4 cellules filles haploïdes de même taille
    \item spermatozoïde = plus petite cellule du corps masculin
\end{itemize}

Méiose chez la femme :
\begin{itemize}
    \item division asymétrique
    \item cellule initiale ~100 µm
    \item M1 : 23 chromosomes par cellule, 90\% cytoplasme dans une cellule et 10\% dans l'autre (globule polaire)
    \item M2 : même principe → second globule polaire
\end{itemize}

Test de paternité : environ 10 microsatellites utilisés.

\section*{12 mars}

Disomie uniparentale : 2 chromosomes mais provenant d'un seul parent (ex : mère donne 2, père 0).\\
Probabilité : 1/10000.

Simplification des probabilités de tailles de microsatellites en les rendant équiprobables.

\section*{20 mars}

\textbf{Épissage (splicing) :} réalisé par les spliceosomes.

\vspace{0.3cm}

Transcription :
\begin{itemize}
    \item ADN polymérase I : 1\%
    \item ADN polymérase III : 99\%
    \item ADN polymérase II (eucaryotes)
\end{itemize}

ADN polymérase II possède un CTD (C-terminal domain), chaîne polypeptidique longue.\\
Elle permet l'ajout des coiffes en 5' des ARN transcrits.

Une protéine se fixe sur cette coiffe.

Chez les procaryotes (E. coli), une seule polymérase.

\vspace{0.3cm}

Exons et introns : éliminés (non codants dans ce contexte).

\section*{23 mars}

Autosomes : 22 paires.\\
Gonosomes : paire X/Y (ex : ornithorynque mâle = 5 X et 5 Y).

Le chromosome Y est souvent exclu car peu de gènes et mutations souvent stérilisantes donc non transmises.

\vspace{0.3cm}

Dominance complète ou incomplète au niveau phénotypique.

Exemple : production de mélanine vs albinisme (dominance complète).

\vspace{0.3cm}

Techniques de microscopie : cryofracture.\\
On congèle un globule puis on le fracture pour observer un plan de rupture.

\vspace{0.3cm}

Anticorps : molécules divalentes liant deux antigènes identiques.

\vspace{0.3cm}

Le phénotype rhésus négatif est récessif.

\vspace{0.3cm}

Dans l'hémoglobine, on enlève la méthionine.




\section*{26 mars}

\textbf{Achondroplasie :}

Tête et corps de taille normale mais bras et jambes plus petits et disproportionnés.

On ne peut pas dire "porteur" car dominant → porteur = atteint d'achondroplasie.

\vspace{0.3cm}

Mutation de novo vs héritée :
\begin{itemize}
    \item De novo : apparaît pour la première fois dans une famille (souvent pour dominant)
    \item Héritée : présente chez un ou les parents
\end{itemize}

90\% des cas sont de novo et 10\% hérités d'un parent atteint.

Homozygote non viable → toujours hétérozygote.

Nanisme thanatophore : naissance possible mais mort peu après.

Allèle FGFR3 muté sur chromosome (fibroblaste).

Plus le père est âgé, plus le risque augmente.

\vspace{0.3cm}

Cellules somatiques : cellules du corps.\\
Cellules germinales : produisent les gamètes.

\vspace{0.3cm}

\textbf{Daltonisme :}

8\% des hommes sont daltoniens rouge/vert.

Récepteurs RGB dans l'œil (opsines) :
\begin{itemize}
    \item Bleu : gène S sur chromosome 7
    \item Vert : gène M sur chromosome X
    \item Rouge : gène L sur chromosome X
\end{itemize}

\section*{27 mars}

Daltonisme récessif chez la femme, mais chez l'homme monogénique (lié à X).

Monogénique = 1 gène.\\
Polygénique = plusieurs gènes (ex : taille).

\vspace{0.3cm}

Hémophilie A et B :
\begin{itemize}
    \item Sur chromosome X
    \item Augmente les risques d'hémorragie
    \item ~30\% de novo
\end{itemize}

Hémarthrose = hémorragie intra-articulaire.\\
Les hémorragies sous-cutanées sont les plus dangereuses.

\vspace{0.3cm}

Mutation PHEX :
rachitisme résistant à la vitamine D (manque de phosphate).\\
Une femme hétérozygote peut être atteinte (jambes arquées).

\vspace{0.3cm}

Gène SRY (différenciation gonades) sur Y.\\
Récepteur testostérone sur X.\\
Femme XY possible → testicules anormaux → risque cancer.

\section*{13 avril}

Compensation de dosage des protéines chez la femme : inactivation d'un X.

Pendant 3 semaines (~19 jours) les deux X sont actifs puis inactivation d'un X.

Corpuscule de Barr :
\begin{itemize}
    \item 0 chez l'homme
    \item 1 chez la femme (ou plus si plusieurs X)
\end{itemize}

Daltonisme rouge-vert :
\begin{itemize}
    \item récessif chez la femme (lié X)
    \item homme porteur direct
\end{itemize}

\vspace{0.3cm}

Histones : lysines chargées positivement.\\
Méthylation → inactivation du X, conservée par mitose.

Méthylation du C uniquement chez mammifères (C suivi de G).

\vspace{0.3cm}

HUMARA (Human Androgen Receptor Assay) : test laboratoire.

Enzymes de restriction :
\begin{itemize}
    \item exo : extrémités
    \item endo : reconnaissance interne
\end{itemize}

Méthylation = inhibition.

Chez phages : ADN méthylé sur A.

\vspace{0.3cm}

\section*{17 avril}

XIC (X Inactivation Center).

Mutation perte ou gain de fonction (pas forcément avantageuse).

Si mutation → perte d'inactivation normale 50-50.

20488 gènes codant pour protéines.

Délétion = perte de fonction.

Désactivation 0-100 possible → femme daltonienne si X muté constant.

Dysplasie ectodermique liée à X.

\vspace{0.3cm}

UPGMA / WPGMA : différences de résultats.\\
Distance en centiMorgan.

\section*{20 avril}

Groupe sanguin et couleur des cheveux sur même chromosome.

Agouti = couleur normale souris.

Ailes vestigiales = mutation.

Plus les gènes sont éloignés, plus crossing-over probable.

CentiMorgan = distance génétique.\\
Valable seulement si gènes sur même chromosome.

50 cM → équivalent non lié (non conclusif).

\section*{27 avril}

Syndrome de Prader-Willi : chromosome 15 paternel.\\
Syndrome d'Angelman : chromosome 15 maternel.

Même région, phénotypes différents.

\section*{30 avril}

Méthylation UBE3A :
\begin{itemize}
    \item maternel = méthylé
    \item paternel = actif
\end{itemize}

Causes Angelman :
\begin{itemize}
    \item 68\% mauvais crossing-over
    \item 13\% mutation ponctuelle UBE3A
    \item 3\% disomie uniparentale paternelle
    \item 6\% défaut empreinte
\end{itemize}

IC = Imprinting Center.

\vspace{0.3cm}

Syndrome généralement sporadique.

\vspace{0.3cm}

Southern blot :
\begin{enumerate}
    \item digestion ADN enzymes de restriction
    \item électrophorèse
    \item fractionnement
\end{enumerate}

Hybridation avec sonde radioactive (aujourd'hui fluorescence).

\section*{1 mai}

Sonde = ADN simple brin.

P32 : demi-vie 15 jours (ancienne méthode).

HindII : insensible à méthylation.\\
HpaII : sensible à méthylation.

IC méthylé maternel, non méthylé paternel.

\vspace{0.3cm}

\section*{11 mai (hors programme)}

Risque cancer peau clair x100 vs peau foncée.

Élastine : matrice extracellulaire, sensible UV, arrêt production après 30 ans.

Carcinome basocellulaire : invasion épithéliale vers derme.

\vspace{0.3cm}

Réparation ADN :
\begin{itemize}
    \item mutation ~1 / $3 \cdot 10^9$ bases
    \item UV, radioactivité, chimio, ROS
\end{itemize}

IRM : pas dommage (champ magnétique).

\vspace{0.3cm}

Sclérodermie pigmentée = "enfants de la lune".

UV → dimères de thymine → déformation ADN.

Apoptose : mort cellulaire programmée (macrophages éliminent).

Après 5 ans : cellules rétiniennes fixées.






































\end{document}